\documentclass[a4paper,10pt,fleqn]{article}
\usepackage[margin=1in]{geometry}
\usepackage{fuzz}

\begin{document}
\title{Vox Hook Context Gaps: A Z Specification}
\author{Punt Labs}
\maketitle

\section{State}

The daemon tracks missing sequence numbers for a single live session. Each
pair denotes an inclusive interval; the implementation stores the pairs in
ascending order. The cutoff bounds what can be deduplicated: a sequence at
or below it cannot be classified as either missing or seen. No gaps after
the highest observed sequence are inferred.

\begin{zed}
  Interval == \nat_1 \cross \nat_1 \\
  Outcome ::= accepted | duplicate | expired
\end{zed}

\begin{axdef}
  maxRanges : \nat_1
\where
  maxRanges = 128
\end{axdef}

\begin{schema}{GapState}
  highest, cutoff : \nat \\
  missing : \finset Interval
\where
  cutoff \leq highest \\
  \# missing \leq maxRanges \\
  \forall lo, hi : \nat_1 | (lo, hi) \in missing @
    cutoff < lo \land lo \leq hi \land hi < highest \\
  \forall a, b, c, d : \nat_1 |
    (a,b) \in missing \land (c,d) \in missing \land (a,b) \neq (c,d) @
    b + 1 < c \lor d + 1 < a
\end{schema}

\begin{schema}{InitGaps}
  GapState
\where
  highest = 0 \\
  cutoff = 0 \\
  missing = \emptyset
\end{schema}

\section{Transitions}

An advancing arrival adds at most one interval. Healing removes the received
number from its enclosing interval, leaving up to two nonempty intervals.
Either can add at most one to the interval count. These operations use
\textit{RetainCandidate} to keep the cutoff unchanged or expire exactly the
oldest interval, including when a heal splits the oldest interval. The
intermediate \textit{candidate} is not externally visible.

\begin{schema}{RetainCandidate}
  \Delta GapState \\
  candidate : \finset Interval
\where
  \# candidate \leq maxRanges + 1 \\
  \# candidate \leq maxRanges \implies
    missing' = candidate \land cutoff' = cutoff \\
  \# candidate = maxRanges + 1 \implies
    (\exists lo, hi : \nat_1 | (lo, hi) \in candidate @ \\
  \t1  (\forall a,b : \nat_1 | (a,b) \in candidate @ lo \leq a) \land
    missing' = candidate \setminus \{(lo, hi)\} \land cutoff' = hi)
\end{schema}

\begin{schema}{AdvanceGaps}
  RetainCandidate \\
  sequence? : \nat_1 \\
  outcome! : Outcome
\where
  sequence? > highest \\
  highest' = sequence? \\
  candidate = missing \cup
    \{ l,h : \nat_1 | l = highest + 1 \land h = sequence? - 1 \land l \leq h \} \\
  outcome! = accepted
\end{schema}

\begin{schema}{HealGap}
  RetainCandidate \\
  sequence? : \nat_1 \\
  outcome! : Outcome
\where
  highest' = highest \\
  \exists lo, hi : \nat_1 |
    (lo,hi) \in missing \land lo \leq sequence? \land sequence? \leq hi @ \\
  \t1 candidate = (missing \setminus \{(lo,hi)\}) \cup \\
  \t2 \{ l,h : \nat_1 | l = lo \land h = sequence? - 1 \land l \leq h \} \cup \\
  \t2 \{ l,h : \nat_1 | l = sequence? + 1 \land h = hi \land l \leq h \} \\
  outcome! = accepted
\end{schema}

\begin{schema}{IgnoreDuplicate}
  \Xi GapState \\
  sequence? : \nat_1 \\
  outcome! : Outcome
\where
  cutoff < sequence? \land sequence? \leq highest \\
  \lnot (\exists lo,hi : \nat_1 |
    (lo,hi) \in missing @ lo \leq sequence? \land sequence? \leq hi) \\
  outcome! = duplicate
\end{schema}

\begin{schema}{RefuseExpired}
  \Xi GapState \\
  sequence? : \nat_1 \\
  outcome! : Outcome
\where
  sequence? \leq cutoff \\
  outcome! = expired
\end{schema}

\section{Implementation checks}

\textbf{Advance and heal:} a jump exposes the skipped intervals; healing
can split an interval without moving the high-water mark.
\textbf{Duplicate framing:} replay cannot alter the gap report or projection.
\textbf{Bounded expiry:} both advance and split may fill the range budget; the
oldest interval expires and an arrival at or below that cutoff changes
nothing. The distinction between duplicate and expired is observable.

\end{document}
